% 03-related-work.tex
\chapter{Stand der Technik und verwandte Arbeiten}
\label{chap:related_work}

In diesem Kapitel wird die relevante Literatur zusammengefasst und diskutiert. Es dient dazu, den aktuellen Forschungsstand aufzuzeigen und die eigene Arbeit in diesen Kontext einzuordnen.

\section{Literaturübersicht}
Verschiedene Studien haben sich bereits mit ähnlichen Fragestellungen beschäftigt. So vergleichen Schmidt und Weber verschiedene Machine-Learning-Verfahren in der Qualitätskontrolle \cite{schmidt2019vergleich}. Andere Arbeiten fokussieren sich auf die Evaluation von Automatisierungsprozessen in der Industrie \cite{mueller2021evaluation}.

\section{Abgrenzung der eigenen Arbeit}
Hier solltest du darstellen, inwiefern sich deine Arbeit von den bisherigen Ansätzen unterscheidet. Welche Forschungslücke wird geschlossen? Wenn du beispielsweise auf offizielle Dokumentationen verweist, kannst du dies ebenfalls tun, wie etwa auf die \LaTeX-Dokumentation \cite{latexproject2023}. 

Oft ist es hilfreich, die Literatur in einer Tabelle gegenüberzustellen (siehe Tabelle~\ref{tab:lit_comparison}).

\begin{table}[htbp]
  \centering
  \caption{Vergleich der untersuchten Literatur und deren Kernaspekte}
  \label{tab:lit_comparison}
  \begin{tabular}{l c l}
    \hline
    \textbf{Quelle} & \textbf{Jahr} & \textbf{Fokus} \\
    \hline
    Schmidt und Weber \cite{schmidt2019vergleich} & 2019 & Machine-Learning-Verfahren \\
    Müller und Richter \cite{mueller2021evaluation} & 2021 & CI/CD-Pipelines \\
    Meyer \cite{meyer2020grundlagen} & 2020 & Grundlagen Softwareentwicklung \\
    \hline
  \end{tabular}
\end{table}

\begin{quote}
  \textit{„Ein Blockzitat kann direkt im Text hervorgehoben werden, um besonders wichtige oder längere Aussagen aus der Literatur wortwörtlich zu betonen.“}
\end{quote}